\section{Especificaciones del host, Docker y observabilidad}
\subsection{Equipo y sistema operativo}
\input{\artifactdir/host.tex}
La captura corresponde al momento de generación del informe. Aunque
se realice en el equipo experimental, no demuestra la frecuencia,
disponibilidad de memoria o carga exacta durante las mediciones.
Las CPU lógicas visibles no deben confundirse con núcleos físicos.

El entorno comunicado utiliza Kubuntu nativo, Docker Engine y Compose,
con cgroups v2. Las versiones comunicadas son MongoDB 9.0.4 y
PostgreSQL 18.6. Una etiqueta de imagen, una versión de biblioteca y la
versión del servidor son evidencias diferentes. Debe cotejarse esa
correspondencia con los artefactos de la cohorte.

\subsection{Arquitectura y servicios}
\input{\artifactdir/compose.tex}
Los servicios e imágenes se extraen de archivos Compose existentes.
Son configuración declarada, no comprobación de salud, persistencia
o límites efectivos. El generador no invoca Docker.

El esquema describe responsabilidades conceptuales y no una topología
de red inferida.
\begin{center}
\begin{tabular}{c}
\fbox{\parbox{0.79\linewidth}{\centering JSONL, manifiesto y equivalencia}}\\[0.2cm]
$\downarrow$\\[0.1cm]
\fbox{\parbox{0.79\linewidth}{\centering Runner: consulta, perfil, concurrencia y ventanas}}\\[0.2cm]
$\downarrow$\\[0.1cm]
\fbox{\parbox{0.34\linewidth}{\centering MongoDB}}\quad
\fbox{\parbox{0.34\linewidth}{\centering PostgreSQL JSONB}}\\[0.2cm]
$\downarrow$\\[0.1cm]
\fbox{\parbox{0.79\linewidth}{\centering CSV, telemetría, preflight y commit}}
\end{tabular}
\end{center}

\subsection{Perfiles de recursos}
\begin{table}[!htbp]\centering
\caption{Perfiles oficiales comunicados.}\label{tab:perfiles}
\begin{tabular}{lrr}\toprule
Perfil & CPU & Memoria (MiB)\\\midrule
\texttt{PROFILE\_A\_BASELINE} & 4 & 4096\\
\texttt{PROFILE\_B\_MEDIUM} & 2 & 2048\\
\texttt{PROFILE\_C\_SEVERE} & 1 & 1024\\
\bottomrule\end{tabular}
\end{table}
Docker y cgroups permiten imponer restricciones \cite{docker,kernel}.
Una cuota equivalente a cuatro CPU no acredita reserva de cuatro
núcleos ni exclusividad. La presencia de cgroups v2 tampoco sustituye
los límites efectivos del cgroup de cada motor.

El límite de memoria no es RSS máximo. La ausencia de swap no equivale
automáticamente al cumplimiento literal de otra opción de swappiness;
las adaptaciones requieren registro.

\subsection{Observabilidad y unidades}
La arquitectura comunicada utiliza Prometheus, cAdvisor y Node Exporter.
La interpretación exige identificar contenedor, etiquetas, intervalo y
ventana. Una serie vacía no equivale a cero. Los contadores requieren
deltas y control de resets.

RSS, working set y uso total del cgroup son señales distintas.
La columna se llama \texttt{max\_rss\_mb}, mientras que la comunicación
reciente usa MiB. Se conservan valores originales, sin convertir ni
renombrar la unidad silenciosamente. La fórmula del recolector y su
correspondencia con la cohorte deben cotejarse.
